\scanpage{11}
étant soumis à une condition de transitivité pour un composé \(v_{\unsure{\circ|o}}u\) de flèches de \(S\) , et à la condition que \(u\unsure{\varepsilon|\in}\underline{M}\) implique que \(\handwritten{\varphi}_u\) soit un isomorphisme .

\handwritten{3}) Définir des foncteurs ``restrictions '' et `` prolongement ''
\[
\operatorname{Res}_X:(S_{\handwritten{/}X})^{\handwritten{\sim}}\to S(X)^{\handwritten{\unsure{\vee|\sim}}}\ ,\quad
\operatorname{Prol}_X:S(X)^{\handwritten{\sim}}\to(S_{\handwritten{/}X})^{\handwritten{\sim}}\ ;
\]
montrer que \(\operatorname{Res}_X\) commute aux limites inductives et aux limites projectives et que \(\operatorname{Prol}_X\) est \uline{pleinement fidèle} , son image essentielle étant formée des faisceaux \(F\) tels que \(\handwritten{\varphi}_u\) soit un isomorp\deleted{g}hisme pour toute flèche \(u\) de \deleted{\unsure{du|\(M\)}} \(S_{\handwritten{/}X}\) .

\handwritten{4}) Définir un isomorphisme d'adjonction faisant de \(\operatorname{Res}_X\) l'adjoint à droite de \(\operatorname{Prol}_X\) . En conclure qu'il existe un morphisme de topos
\[
f:S(X)^{\handwritten{\sim}}\to(S_{\handwritten{/}X})^{\handwritten{\sim}}
\]
tel que \(f^{\handwritten{*}}=\operatorname{Res}_X\) et que le foncteur \(f^{\handwritten{*}}\) admet un adjoint à gauche \(f_!=\operatorname{Prol}_X\) . Montrer que \(\operatorname{Prol}_X\) transforme les faisceaux abéliens en faisceaux abéliens . Montrer que \(f:S(X)^{\handwritten{\sim}}\handwritten{\longrightarrow}(S_{\handwritten{/}X})^{\handwritten{\sim}}\) est un \uline{plongement} de topos (\handwritten{9}.\handwritten{1}.\handwritten{1}) et que par suite ``le petit topos '' \(\handwritten{S(X)^{\sim}}\) est \deleted{\unsure{un sous-topos (9.1}} équivalent à un sous-topos du ``gros topos '' \((S_{\handwritten{/}X})^{\handwritten{\sim}}=S^{\handwritten{\sim}}_{\handwritten{/}X}\) .

\handwritten{5}) Soit \(F\) un faisceau abélien flasque de \(S^{\handwritten{\sim}}_{\handwritten{/}X}\) (V \handwritten{4}) . M\textsuperscript{O}ntrer que \(\operatorname{Res}_X F\) est \(S(X)\)-acyclique\inserted{\handwritten{(V 4.1)}} . En déduire que pour tout faisceau abélien \(F\) de \(S^{\handwritten{\sim}}_{\handwritten{/}X}\) et tout entier \(q\) , le groupe \(H^q(X,F)\) (cohomologie calculée dans \(S^{\handwritten{\sim}}_{\handwritten{/}X}\) ) est canoniquement isomorphe au groupe \(H^q(X,\operatorname{Res}_X F)\) (cohomologie calculée dans \(S(X)^{\handwritten{\sim}}\) ) .

\handwritten{6}) Déduire de 5) que pour tout faisceau abélien \(G\) \deleted{sur} de \(S(X)^{\handwritten{\sim}}\) et tout entier \(q\) le groupe \(H^q(X,G)\) \deleted{(resp.} (cohomologie calculée dans \(S(X)^{\handwritten{\sim}}\) ) est canoniquement isomorphe aux groupes \(H^q(X,f_{\handwritten{*}}G)\) et \(H^q(X,\operatorname{Prol}_X G)\) (cohomologie calculée dans \(S^{\handwritten{\sim}}_{\handwritten{/}X}\) ) .
% archive sheet number: 10
